% ================================================================
% chapters/parts/ch07-c.tex —— 第 7 章 原位红外光谱方法（分片 c）
% 覆盖：原书 §7.5（加氢精制催化剂活性相与助剂作用，硫化物催化剂）
%        §7.6（原位红外应用于反应机理研究，含光催化水氧化）
%        以及 §7.6 之后的「电化学／光电催化原位红外」部分
%        书内 p.489—518 / PDF p.508—537
% 注：源文本前 93 行（书内 p.489 上的图 7-77、图 7-78，甲烷在 Zn-ZSM-5 上的
%     解离吸附）属 §7.4 的收尾，归分片 b 的范围，本分片不重复。
% 版式自查（WRITER_BRIEF §7.3，单栏版心 165 mm；\onecolwidth ≈ 132 mm）：
%   图 \ref{fig:ch07c-transient}：groupplot 两列，2×54 + 8 = 116 mm，
%     加左侧轴标签与刻度约 13 mm → ≈129 mm ≤ 165 mm  ✓
%   表 \ref{tab:ch07c-probe}、\ref{tab:ch07c-criteria}、\ref{tab:ch07c-formate}
%     均用 tabularx 定宽 \onecolwidth，无「定宽 p{} 凑总宽」。不用 xltabular。
%   declare function 的参数只用 \p（避开 \x、\u），函数体乘法全部加括号；
%     groupplot 不给公共 domain，逐面板显式给 domain，避免 Dimension too large。
% ================================================================

\subsection{加氢精制催化剂活性相与助剂作用（硫化物催化剂）}
\label{sec:ch07c-sulfide}

加氢精制是炼油工业的四大支柱工艺之一。其催化剂牌号有上百种，但大部分是
Co、Mo、Ni、W 的不同组合，工业催化剂载体大多为 \ce{Al2O3}，而使用状态
\emph{均为硫化物}\src{S1 §7.5，书内 p.489 / PDF p.508}。围绕活性相与助剂
作用，文献提出了多个模型（如 CoMoS 相模型与协同作用模型），长期争论不休；
本书作者点出的主要原因是\textbf{原位信息不多}\src{S1 §7.5.4，书内 p.493 /
PDF p.512}。这正是原位红外在这个体系上的位置：它不是给出「催化剂是什么」，
而是在工作状态下给出「有几类可被探针分子接近的中心、助剂改变了哪一类」。

\begin{probout}
本节要输出三件事：①表面上存在哪几类\textbf{配位不饱和中心}（coordinatively
unsaturated site，CUS），各自是什么价态与配位状态；②助剂（Co、Ni、Ru）
改变的是\emph{哪一类}中心的数目或电子状态；③这些改变怎样与加氢脱硫（HDS）、
加氢脱氢（HDY）活性对应起来。要注意的是，红外能可靠给出的是「中心类型」
与「中心的相对变化」；中心数目的绝对值必须由滴定实验（化学吸附或探针分子
饱和吸附）标定后才能谈，否则强度只能读成相对量。
\end{probout}

\subsubsection{硫化物表面在红外测量中的三个特殊问题}

硫化物催化剂的谱学行为与氧化物催化剂有实质差别，读谱前必须先处理这三件事，
否则得到的是「强度失真后的假变化」。

\paragraph{（1）全谱区吸收与散射使基线不稳。}硫化态的 Co--Mo 或 Ni--Mo
催化剂外观为黑色，在可见—红外区有很强的电子吸收；样品的散射又随粒径、
压片或装样方式改变。透射测量中，样品自身吸收与散射把光耗掉，谱带信噪比低、
基线随波数漂移；漫反射测量则要先做 Kubelka--Munk 变换才能把反射率变成与
浓度近似成线性的量
\begin{equation}
  F(R_{\infty})=\frac{(1-R_{\infty})^{2}}{2R_{\infty}}=\frac{K}{S},
  \label{eq:ch07c-km}
\end{equation}
式中 $R_{\infty}$ 为半无限厚样品的漫反射比，$K$ 为吸收系数，$S$ 为散射系数。
\textbf{适用条件是变换成立的前提}：样品必须足够厚（光未穿透）、且在所比较的
谱段内 $S$ 近似不变。硫化物样品粒径分布宽，硫化度改变时表面硫覆盖度也改变，
$S$ 并非常数，故 $F(R_{\infty})$ 的变化里混有散射项。可行的做法是在同一张谱
内取带强度比（例如 $A_{2072}/A_{1694}$），不跨样品比绝对强度。
\tocheck{未对照原页面图像：硫化物自吸收与散射的这段讨论源文未展开，本段按
红外反射光谱的通用处理（Kubelka--Munk 变换及其适用条件）补充，
式~\eqref{eq:ch07c-km} 为学科标准形式，未与原书符号约定逐字核对。}

\paragraph{（2）自吸收与强吸收带附近的谱带失真。}当吸收系数在某个波数附近
出现极大（强吸收带、Reststrahlen 带与 Christiansen 效应所在的波数区间），
颗粒内部的折射与散射会同时变化，谱带的位置和形状都会畸变：带可能变成
不对称，甚至出现表观位移。硫化物在低波数区（骨架与 \ce{Mo-S} 相关振动）
吸收很强，恰是这类畸变最严重的地方。因此报告低波数带位置时，必须同时给出
样品稀释比例与装样方式，并且用同一批装样的谱做相互比较。

\paragraph{（3）羟基与硫醇基必须区分。}硫化后的表面同时存在残余
O--H 与 S--H 两类基团：O--H 伸缩在 $3400$—$3700\;\mathrm{cm^{-1}}$
（游离羟基更高、氢键羟基更低，且随脱水前处理位移），S--H 伸缩在
$2550$—$2600\;\mathrm{cm^{-1}}$，位置分得很开；真正的困难在于
\textbf{S--H 带很弱}——S--H 的偶极变化小，消光系数远低于 O--H，容易被噪声
或水汽、\ce{CO2} 的干扰带淹没。判据是两条：①用 D 交换（\ce{D2} 或
\ce{D2S} 处理）复测，\ce{X-H} 带应按约 $1.3$—$1.4$ 倍的位移比移向低波数
（见式~\eqref{eq:ch07c-isotope} 与其后的说明）；②在惰性气氛下把样品脱水到
同一程度再比较，排除羟基随含水量漂移带来的假变化。把 $3400$—
$3700\;\mathrm{cm^{-1}}$ 的带直接当成「酸性羟基」、把 $2550$—
$2600\;\mathrm{cm^{-1}}$ 的弱带忽略掉，是这一体系上最常见的两个起点错误。

\paragraph{（4）硫化度是变量，不是条件。}同一种催化剂在不同硫化条件下，
表面 CUS 的\textbf{类型分布}会变，谱带因此改变。源文给出的关键实验是：
用漫反射红外池、以 $5\%$ \ce{H2S}/\ce{H2} 在常压和高压（$5.1$\,MPa）下
硫化同一种 Co--Mo/\ce{Al2O3}，再吸附 NO 作探针。结果是 NO 吸附在 Mo 中心
的 $1690\;\mathrm{cm^{-1}}$ 带随硫化压力升高而明显减弱、直至消失；而在机械
混合的 Co--Mo/\ce{Al2O3} 上，这一压力效应不明显\src{S1 §7.5.4，书内 p.495—
496 / PDF p.514—515}。也就是说，
\textbf{随硫化压力增加，Co 的 CUS 增加而 Mo 的 CUS 明显减弱}——读一张硫化态
谱图时若不写明硫化压力，谱带强弱就无法互比。硫化度的定量表述通常用
S/(Co$+$Mo) 原子比（由 XPS 给出）或探针分子滴定的 CUS 数密度。
\tocheck{未对照原页面图像：原书 §7.5.4 只给硫化压力（常压与 $5.1$\,MPa）与
谱带变化，未给「硫化度」的定量定义式；此处按学科通用定义（S/(Co$+$Mo)
原子比）补充，未与原书符号约定核对。}

\begin{measure}
这一体系的测量条件本身就是结论的一部分，报告时必须写全：
\begin{itemize}
  \item \textbf{硫化}：\ce{H2S}/\ce{H2}（源文用 $5\%$）常压或加压硫化，
        硫化压力与硫化温度都要给——它们决定 CUS 的类型分布（见上文第 4 点）。
  \item \textbf{探针分子}：CO 室温吸附，分压要低。气相 CO 本身的伸缩带落在
        $2143\;\mathrm{cm^{-1}}$，分压偏高时会与表面线式、桥式吸附带叠在
        同一区间，把「吸附态的变化」读成「气相的贡献」；NO 与 CO 共吸附时，
        源文明确指出共吸附态的形成与\textbf{吸附次序}密切相关，次序必须写清
        \src{S1 §7.5.3，书内 p.492 / PDF p.511}。
  \item \textbf{样品状态}：硫化物样品不能在空气中装样与长时间暴露（再氧化会
        同时改变表面硫覆盖度与 Mo 的价态）；先按与红外实验相对应的条件脱水，
        再原位采集谱图\src{S1 §7.5.1，书内 p.491 / PDF p.510}。
  \item \textbf{定量纪律}：只用同一张谱内的带强度比；报告硫化压力、硫化温度、
        探针分压、吸附温度与样品稀释比例。
\end{itemize}
\end{measure}

\subsubsection{氧化物前驱体：\ce{MoO3/Al2O3} 与 \ce{WO3/Al2O3} 的还原态判读}

\begin{principle}
双探针法的原理是\textbf{化学计量取代关系}：CO 与 NO 都只吸附在配位不饱和的
金属中心上，两者对同一中心的亲和力不同。单独吸附 CO 时，不同中心上的 CO
带常叠在一起；引入 NO 后，被 NO 夺占的那一类中心上的 CO 带消失或位移，
于是\textbf{共吸附谱相当于给单吸附谱做了一次「分峰对照实验」}。
定量仍以朗伯—比尔定律为基准：$A=\varepsilon c l$（$A$ 为积分吸光度、
$\varepsilon$ 为积分消光系数、$c$ 为单位面积上的表面浓度、$l$ 为有效光程），
把「一位一分子」的化学计量滴定近似代入，$c$ 才是该类型中心的面密度。
\end{principle}

\textbf{\ce{MoO3/Al2O3}（$10\%$ Mo）。}把原位激光拉曼与 CO、NO 双探针红外
联用，可以追踪不同还原温度下的价态与配位分布\src{S1 §7.5.1，书内 p.490—
491 / PDF p.509—510}：

\begin{itemize}
  \item $473$\,K 还原后\textbf{不出现}任何 CO 吸附峰——表面没有配位不饱和的
        中心，这一步是后面所有归属的零点。
  \item $573$\,K 以上还原后，CO 吸附出现约 $2183$ 与 $2174\;\mathrm{cm^{-1}}$
        的带，且峰强随还原温度升高而增大。
  \item CO 与 NO 共吸附后，CO 带分为两个（约 $2200\;\mathrm{cm^{-1}}$ 与
        另一带，后者数值在 OCR 中不可辨）：说明单独吸附时看到的是两种中心的
        叠合，共吸附时其中一种 CO 受 NO 影响红移，两种中心由此被分开。
  \item NO 在另一类 Mo 中心上呈孪生吸附，带位于 $1812$ 与
        $1708\;\mathrm{cm^{-1}}$。
  \item 拉曼回查：氧化态样品表面存在三种 Mo 物种（多钼氧物种、单钼氧物种，
        以及载体金属离子杂质形成的钼酸盐）；还原后出现 $280$ 与
        $748\;\mathrm{cm^{-1}}$ 新峰，归属为还原态 Mo 物种，主要由五配位
        \ce{Mo^{5+}} 与四配位 \ce{Mo^{4+}} 组成。
\end{itemize}

由配位数可以理解红外结果：\ce{Mo^{5+}} 只有 1 个配位不饱和位，只能吸附
1 个 CO（带位约 $2200\;\mathrm{cm^{-1}}$）；\ce{Mo^{4+}} 有 2 个位，可形成
孪生态并可再吸附 NO，共吸附后出现 $2134\;\mathrm{cm^{-1}}$ 一带。这条链条
还有独立的毒化证据：在氢解与同位素反应中，极少量的 CO 与 NO 就能完全毒化
催化剂，说明这些反应的活性与\textbf{可配位的 \ce{Mo^{4+}}} 有关——只有
\ce{Mo^{4+}} 既与 CO 作用又与 NO 作用。
\tocheck{未对照原页面图像：$2183/2174$、$1812/1708$、$2134\;\mathrm{cm^{-1}}$
四组数值在源文本中均可辨，但与它们配对的第二组带位（共吸附后分裂出的另一带、
孪生峰的第二个带位）在 OCR 中丢失；还原温度一处作 $600$\,K、另一处作
「$600$ 以下不出现新拉曼峰」，本分片按「约 $600$\,K 出现」表述，未核对。}

\textbf{\ce{WO3/Al2O3}（$23.4\%$ W）。}\ce{WO3/Al2O3} 很难还原，表面活性相
的表征信息因而很少\src{S1 §7.5.3，书内 p.492—493 / PDF p.511—512}。同样的
双探针做法给出：\ce{H2} 还原后随还原温度升高，CO 吸附逐步出现 $2198$、
$2176$ 等带；CO 与 NO 共吸附时 $2176$ 与 $2154\;\mathrm{cm^{-1}}$ 带消失，
在 $2117\;\mathrm{cm^{-1}}$ 出现新带，同时 NO 的 $1801$、$1777$、
$1691\;\mathrm{cm^{-1}}$ 带出现；真空脱附实验显示 $2198\;\mathrm{cm^{-1}}$
带不稳定。据此可区分出三种 W 中心：\ce{W^{5+}} 与配位状态不同的两种
\ce{W^{4+}}；共吸附时形成 \ce{(CO)W^{4+}(NO)2} 型共吸附态（新带
$2117\;\mathrm{cm^{-1}}$）。共吸附态的形成与 CO、NO 的吸附次序密切有关，
载量与还原程度也强烈影响其相对量。
\tocheck{未对照原页面图像：\ce{WO3/Al2O3} 一段中 $1801/1777/1691$ 三个 NO
带位可辨，$2117\;\mathrm{cm^{-1}}$ 之后还有一个带位与「三空穴」类表述在 OCR
中损坏；本段只保留了可辨数值，未补全缺失数值。}

\subsubsection{硫化态活性相与助剂作用：一个读图实例}

\begin{dataflow}
读图对象是硫化 \ce{Co/Al2O3}、\ce{Mo/Al2O3}、\ce{Co-Mo/Al2O3} 以及二者
机械混合样的 CO／NO 共吸附谱（原书图 7-83～图 7-85）\src{S1 §7.5.4，
书内 p.493—495 / PDF p.512—514}。判读链条是三步：单吸附谱给出「有哪几种
CO 吸附态」→ 共吸附谱给出「哪几种态能被 NO 取代」→ 与机械混合样对照，
给出「助剂是通过形成相互作用相起作用，还是通过溢流起作用」。三步缺一不可：
只有单吸附谱，无法把叠合的带分开；只有共吸附谱，没有单吸附谱的对照，
位移无从谈起；没有机械混合样，就无法把「相互作用相」与「溢流」这两种机制
分开。谱带、共吸附变化与归属判读的汇总见 \tabref{tab:ch07c-probe}。
\end{dataflow}

\begin{obs}
（只列谱学事实，不含解释）
\begin{itemize}
  \item[(a)] 硫化 \ce{Co/Al2O3} 上 CO 呈线式（$2057\;\mathrm{cm^{-1}}$）与
        桥式吸附（桥式带位在 OCR 中不可辨）；硫化 \ce{Mo/Al2O3} 上只有线式
        带（$2112\;\mathrm{cm^{-1}}$，同书另一处作 $2097\;\mathrm{cm^{-1}}$）；
        硫化 \ce{Co-Mo/Al2O3} 上 CO 带在 $2065\;\mathrm{cm^{-1}}$。
  \item[(b)] CO 与 NO 共吸附后：\ce{Co/Al2O3} 上线式与桥式 CO 带都消失，
        出现吸附态 NO 的 $1557$ 与 $1791\;\mathrm{cm^{-1}}$ 带；
        \ce{Mo/Al2O3} 上 $2097\;\mathrm{cm^{-1}}$ 带强度没有明显变化，
        同时出现 NO 的 $1790\;\mathrm{cm^{-1}}$ 带；\ce{Co-Mo/Al2O3} 上 CO
        带由 $2065$ 移到 $2072\;\mathrm{cm^{-1}}$ 且强度明显下降，Co 中心上的
        NO（$1790\;\mathrm{cm^{-1}}$）与 Mo 中心上的 NO（$1694\;
        \mathrm{cm^{-1}}$）同时出现。
  \item[(c)] 机械混合的 $6\%$ \ce{Co/Al2O3} $+$ $10\%$ \ce{Mo/Al2O3}
        （或两者分别预硫化后再混合，此时 Co--Mo 之间的相互作用可以忽略）
        的共吸附结果，与硫化 \ce{Co-Mo/Al2O3} 的结果相近。
  \item[(d)] 用 NO 作探针时，NO 吸附在 \ce{Mo/Al2O3}、\ce{Co/Al2O3}、
        \ce{Co-Mo/Al2O3} 以及机械混合样上各出现两个带，但带位彼此不同；
        Co 的存在减少了表面上可吸附 NO 的 Mo 中心数。
  \item[(e)] 硫化压力由常压升到 $5.1$\,MPa 时，NO 在 Mo 中心的
        $1690\;\mathrm{cm^{-1}}$ 带明显减弱直至消失；机械混合样上这一效应
        不明显。
\end{itemize}
\end{obs}

\begin{infer}
（每条推断都对应上面的事实编号）
\begin{itemize}
  \item 由 (b) 得：\textbf{NO 能取代 Co 中心上的 CO，不能取代 Mo 中心上的
        CO}。因此 $1790$ 与 $1694\;\mathrm{cm^{-1}}$ 可分别当作 Co 中心与
        Mo 中心的标记带，取代关系本身就是「分峰」的依据。
  \item 由 (a) 与 (b) 得：硫化 \ce{Co-Mo/Al2O3} 的 CO 谱带是 Co 中心与
        Mo 中心吸附 CO 的\textbf{加和}；其中 Mo 中心的吸附带相对
        \ce{Mo/Al2O3} 位移了 $25\;\mathrm{cm^{-1}}$（$2097\to2072$）。
        按「电子密度升高、$v_{\mathrm{CO}}$ 红移」的判读，这说明
        \ce{Co-Mo/Al2O3} 中 Mo 中心的\textbf{还原度更高}。
  \item 由 (c) 得：由于机械混合时 Co 与 Mo 之间不存在相互作用相，结果却与
        \ce{Co-Mo/Al2O3} 相近，\textbf{助剂作用不必以形成 Co--Mo 相互作用相
        为前提}；用「\ce{H2} 在 Co 中心解离、氢溢流到邻近 Mo 中心将其还原」
        即可解释 Mo 中心的调变。这条推断与另一组观测自洽：$2\%$ Co 的加入
        并未改变 Mo 的可还原性（还原后同样只有很弱的还原态 Mo 峰），
        但 Co--Mo--O 相互作用相的生成减少了表面 Mo 氧物种浓度，使还原态
        Mo 物种总量减少、CO／NO 吸附峰随 Co 含量增加而逐渐降低
        \src{S1 §7.5.2，书内 p.491—492 / PDF p.510—511}。
  \item 由 (d) 得：NO 可以\textbf{滴定}表面 Co 中心与 Mo 中心的数目——
        这正是「强度给出数目」的用法，但前提是先用标样或饱和吸附标定
        每个中心的 NO 化学计量数。
  \item 由 (e) 得：硫化压力改变的是 CUS 的\textbf{形成过程}，而不是探针分子
        的吸附性质；该推断只在「被比较的谱取自同一装样、同一探针条件」时成立。
\end{itemize}
\end{infer}

\begin{table}[htbp]
\centering
\caption{硫化态 Co--Mo/\ce{Al2O3} 体系上 CO 与 NO 双探针共吸附的谱带归属与判读
\src{据 S1 §7.5.4 及原书图 7-83～图 7-87，书内 p.494—497 / PDF p.513—516 整理。}}
\label{tab:ch07c-probe}
\begin{tabularx}{\onecolwidth}{@{}>{\raggedright\arraybackslash}p{25mm}
  >{\raggedright\arraybackslash}p{29mm} X@{}}
\toprule
样品 & CO 吸附带 $/\mathrm{cm^{-1}}$ & 共吸附 CO$+$NO 后的变化与判读 \\
\midrule
硫化 \ce{Co/Al2O3} & 线式 2057；桥式（带位未辨） &
线式与桥式带均消失，出现吸附态 NO 的 1557 与 1791：该中心上的 CO 能被
NO 取代，记为 Co 中心 \\
硫化 \ce{Mo/Al2O3} & 线式 2112（另处作 2097） &
2097 带强度无明显变化，同时出现 NO 1790：Mo 中心上的 CO 不能被 NO 取代 \\
硫化 \ce{Co-Mo/Al2O3} & 2065 &
CO 带移到 2072 且强度明显下降；Co 中心 NO 1790 与 Mo 中心 NO 1694 同时
出现：谱带是两类中心吸附 CO 的加和，Mo 中心位移 $25\;\mathrm{cm^{-1}}$ \\
机械混合 $6\%$ \ce{Co/Al2O3} $+$ $10\%$ \ce{Mo/Al2O3}
（或分别预硫化后混合） & 与硫化 \ce{Co-Mo/Al2O3} 相近 &
结果与上一行相近：Co 与 Mo 之间无须相互作用相，氢溢流即可调变 Mo 中心 \\
常压硫化 $\to$ $5.1$\,MPa 硫化 & NO 在 Mo 中心的 1690 带减弱至消失 &
硫化压力改变 CUS 的形成：Co 的 CUS 增加、Mo 的 CUS 明显减弱；
机械混合样上该效应不明显 \\
\bottomrule
\end{tabularx}
\end{table}

\tocheck{未对照原页面图像：表中 \ce{Mo/Al2O3} 的 CO 线式带位在源文本中出现
两个值（$2112$ 与 $2097\;\mathrm{cm^{-1}}$），本表并记两值；
\ce{Co/Al2O3} 的桥式 CO 带位、以及 (b) 中共吸附分裂出的第二个 CO 带位在
OCR 中丢失。}

\textbf{助剂不止 Co。}用同一套双探针法研究硫化态 Ru--Mo/\ce{Al2O3} 与
Ru--\ce{Al2O3} 时，Ru 的加入大大提高了 CO 与 NO 吸附峰的强度，说明 Ru 的
存在产生了更多配位不饱和的 Mo 中心，由此提高了催化剂的 HDS 与 HDY 活性
\src{S1 §7.5.4，书内 p.495 / PDF p.514}。探针分子的选择也不限于 CO／NO：
氮化铁（\ce{Fe2N}，负载于 fumed \ce{SiO2} 上）上 CO 吸附给出
$2040\;\mathrm{cm^{-1}}$ 的线式带（Fe 位点）；与 \ce{Fe/SiO2} 相比，\ce{Fe2N}
样品上 \ce{CO2} 与 \ce{NH3} 类物种出现 $1610$、$1550\;\mathrm{cm^{-1}}$
新带，其中 $1610\;\mathrm{cm^{-1}}$ 归属为 \ce{NH2} 的变形振动，而
$2208\;\mathrm{cm^{-1}}$ 归属为 N 位点上的 CO 形成 \ce{NCO} 吸附物种——
这表明除 Fe 位点之外，表面 N 位点可能具有高反应活性；再用 \ce{NH3} 与
\ce{H2} 作探针的原位红外进一步支持了「N 位点为高活性反应位点」这一结论
\src{S1 §7.5.4，书内 p.497—498 / PDF p.516—517}。

\begin{pitfall}
\textbf{把「探针分子能吸附」直接读成「活性位数目」。}吸附量给出的是
\emph{可接近}的中心数，它随探针分子的大小、分压、吸附温度而变；强度要换算成
数目，必须先有化学计量数与消光系数的标定。$A_{2072}/A_{1694}$ 这类比值只能
说明两类中心的相对变化，而且还要假设两类中心上 CO 的积分消光系数相同。

\textbf{只用带位归属，不做对照实验。}本书的做法之所以站得住，是因为每一步都有
对照：单吸附对共吸附（分峰）、\ce{Co-Mo} 对机械混合（区分相互作用相与溢流）、
常压硫化对高压硫化（区分制备变量）。只有一张谱上的一个带位，位置再「符合
文献区间」，也不能确定中心类型。

\textbf{忽略气相物种的贡献。}气相 CO 的 $2143\;\mathrm{cm^{-1}}$ 带、
气相 \ce{CO2} 的 $2349\;\mathrm{cm^{-1}}$ 带都会落进关心区间；探针分压偏高时，
「吸附带位移」可能只是气相带叠加后的表观位移。读图前先确认气相带已被扣除
（背景谱、吹扫或低压吸附）。

\textbf{用氧化物前驱体的结论代替硫化态结论。}硫化后 Mo 的价态、配位环境与
表面硫覆盖度都变了，\ce{MoO3/Al2O3} 上「$473$\,K 还原后无吸附峰」这类结论
不能直接搬到硫化态。反过来，硫化态上测到的 CUS 分布也不能用来推断前驱体的
表面结构。

\textbf{忽略原位池窗口与样品自吸收造成的强度失真。}窗片材料的可用波数范围、
样品稀释比例与装样厚度都会改变基线；跨样品、跨批次直接比强度，最容易把
「装样不同」读成「活性相不同」。
\end{pitfall}

\begin{crosscheck}
\textbf{与拉曼互查物相与聚合度。}拉曼给出 \ce{Mo{=}O} 端基（$900$—
$1000\;\mathrm{cm^{-1}}$）、\ce{CoMoO4}（$932\;\mathrm{cm^{-1}}$）以及还原态
Mo 物种的低波数带；红外给出探针分子可接近的中心。两者一起才能说清
「表面上有什么相」与「这些相上有几个可用的位」\src{S1 §7.5.1—§7.5.2，
书内 p.490—492 / PDF p.509—511}。

\textbf{与 XPS／EPR 分工。}XPS 与 EPR 能给出不同价态离子的总相对丰度，但
\textbf{不能直接对应真正的活性中心}（源文明确指出这一点）——因为体相与表面
不可接近的离子同样计入丰度。红外探针给的恰是「配位不饱和、能被分子接近」
的那一部分，两者之差本身就是信息。

\textbf{与活性数据回查。}谱带变化只有与 HDS／HDY 活性同向变化才有解释力；
若谱带变了而活性不变，优先怀疑归属或强度失真。
\end{crosscheck}

\subsection{原位红外光谱应用于反应机理的研究}
\label{sec:ch07c-mech}

前面各节讨论的吸附数据，绝大多数是在\textbf{反应没有发生时}测得的。而要阐明
催化作用机理，仅有反应物与产物分别测得的吸附数据是不够的：反应条件下（或
反应定态下）吸附物种的类型、结构、性能，与吸附条件下可能有很大差别。更
麻烦的是，反应条件下的表面不只存在一种吸附物种，而且\emph{不是所有吸附物种
都一定参与反应}，因此如何在多种吸附物种中识别出参与反应的「中间物」，是
这一节的核心课题\src{S1 §7.6，书内 p.498 / PDF p.517}。

\begin{probout}
要输出的是：某一表面物种\textbf{是不是}反应中间体，以及它在反应路径上的位置
（生成它的步骤、消耗它的步骤）。难点在于「峰的存在」与「峰的参与」是两件事：
一张定态谱上的强带，可能只是储量很大的旁观者；一个很弱、甚至只在瞬态出现的
带，才可能是真正的中间体。因此本节不提供「某峰就是某中间体」的对照表，
而是给出一组\textbf{可执行、可否证的判据}，以及每条判据在本书实例中的落点。
\end{probout}

\subsubsection{四条判据：怎样把「吸附态」与「中间体」分开}

\begin{principle}
\textbf{判据一：序列上的消长关系。}中间体必须\textbf{在时间或温度序列上}
表现出「先生成、后消耗」的消长：其带强度应在产物快速上升之前达到极大，
且其消耗速率与产物生成速率同量级。只给一张定态谱，位置再合理也无法判定
中间体；只有一条单调上升或单调下降的曲线，同样不能判定——单向变化既可能是
中间体的积累，也可能只是吸附平衡随温度的移动。判据一的可视化见
\figref{fig:ch07c-transient}(a)：候选中间体的极大值出现在产物速率快速上升
\emph{之前}，而储量大的甲酸盐只是缓慢下降，两者与产物速率的关联强度完全不同。

\textbf{判据二：速率一致性（稳态动力学）。}在循环系统中用压力与组成变化测
总包反应速率，同时用红外测同一温度、同一覆盖度下该表面物种的分解速率：
\begin{itemize}
  \item 若表面物种的分解速率\emph{等于}总包速率，则该分解就是控速步骤，
        该物种是中间体；
  \item 若表面物种的分解速率\emph{远小于}总包速率，它只能是旁观者。
\end{itemize}
本书给出的定量例正是后者与前者各一：真空、反应温度下 \ce{Al2O3} 上
\ce{HCOO-} 吸附态的分解速率比同温同覆盖度下甲酸蒸气分解速率\textbf{小两个
数量级}；而在 \ce{ZnO} 上，\ce{HCOO-} 的分解速率\textbf{等于}甲酸分解的
总包速率\src{S1 §7.6.1，书内 p.499—501 / PDF p.518—520}。同一个表面物种、
同一个反应，在两个氧化物上一个不是中间体、一个是中间体——这一对例子说明
「中间体」不是物种的固有属性。

\textbf{判据三：同位素置换。}把带归属从「位置判断」变成「比值判断」的依据是
谐振近似
\begin{equation}
  \tilde{\nu}=\frac{1}{2\pi c}\sqrt{\frac{k}{\mu}},\qquad
  \mu=\frac{m_{1}m_{2}}{m_{1}+m_{2}},\qquad
  \frac{\tilde{\nu}(\mathrm{XD})}{\tilde{\nu}(\mathrm{XH})}
  =\sqrt{\frac{\mu_{\mathrm{XH}}}{\mu_{\mathrm{XD}}}},
  \label{eq:ch07c-isotope}
\end{equation}
式中 $\tilde{\nu}$ 为波数，$c$ 为光速，$k$ 为键力常数，$\mu$ 为约化质量。
当 X 比 H 重得多时，$\mu_{\mathrm{XH}}\approx m_{\mathrm{H}}$、
$\mu_{\mathrm{XD}}\approx 2m_{\mathrm{H}}$，故\ce{X-H} 换成 \ce{X-D} 后频率
降到约 $1/\sqrt{2}\approx0.71$ 倍，即 H 物种的波数是 D 物种的约 $1.4$ 倍。
\textbf{适用条件与实测偏差}：这是谐振近似；真实观测到的位移比通常在
$1.3$—$1.4$ 之间（非谐性与模式耦合使比值小于 $1.414$）。判据的用法不是
「必须等于 $1.414$」，而是\emph{同一张谱上若有一个带按这个比值移动，就支持
它含 \ce{X-H}}；若某个自称是 \ce{X-H} 的带在 D 交换后完全不动，归属就有问题。
置换实验还有更强的一种用法：先把反应在 D 标记的反应物上做到定态，再迅速换回
H 标记的反应物，同时跟踪表面物种与\textbf{气相组成}——若表面 D 物种消失时，
D 是以\emph{反应物分子}的形式回到气相，而不是以含 D 的产物出现，说明它主要
走可逆脱附通道，不是通往产物的中间体（对应判据四的实例）。

\textbf{判据四：置换、毒化与共吸附组合。}三条独立的操作性检验：
①\textbf{选择性毒化或预处理}：用一个只减少候选物种、不改变其他条件的处理，
看活性是否同比例下降——量减而活性不变，即可否掉该物种。②\textbf{组分来源
组合}：分别通入不同反应物组合，看该物种能否生成；只有真实反应物组合才能
生成它，才支持它处在主路径上。③\textbf{与产物生成温度的对应}：程序升温表面
反应（TPSR）中，该物种分解/脱出的温度区间应与产物出现的温度区间对应。
\end{principle}

四条判据的操作方式、支持「中间体」的结果，以及本书的实例与反例，
汇总于 \tabref{tab:ch07c-criteria}。

\begin{table}[htbp]
\centering
\caption{判定「反应中间体」的四条判据、操作方式与本书实例（含反例）
\src{据 S1 §7.6.1，书内 p.498—501 / PDF p.517—520 归纳。}}
\label{tab:ch07c-criteria}
\begin{tabularx}{\onecolwidth}{@{}>{\raggedright\arraybackslash}p{19mm} X
  >{\raggedright\arraybackslash}p{31mm} X@{}}
\toprule
判据 & 怎么做 & 支持「中间体」的结果 & 本书实例与反例 \\
\midrule
一、消长关系 & 在同一升温或时间序列上同时记录候选物种与产物信号，比较
极大值与拐点的先后 & 候选物种的极大值出现在产物快速上升\textbf{之前}，
其后被消耗 & \ce{Al2O3} 上甲酸盐在升温中缓降而与产物速率不同步，
不能据此判为中间体（\figref{fig:ch07c-transient}(a) 为模型示意） \\
二、速率一致性 & 循环系统测总包速率；同时用红外测同一温度、同一覆盖度下
该表面物种的分解速率 & 表面物种分解速率 $=$ 总包速率 &
\ce{ZnO} 上 \ce{HCOO-} 分解速率 $=$ 总包速率（是中间体）；
\ce{Al2O3} 上 \ce{HCOO-} 的分解速率比甲酸蒸气分解速率小两个数量级
（不是中间体） \\
三、同位素 & D 交换后复测，或把 D 标记反应物做到定态再换回 H 标记反应物，
同时跟踪表面与气相 & \ce{X-H} 带按约 $1.3$—$1.4$ 倍比值移动；
表面 D 物种消失时 D 以\emph{产物}形式出现 & \ce{Al2O3} 上 DCOOD 的
O--H 带由 3580 移到约 2650、C--H 带移到约 2200、$\delta(\mathrm{C{-}H})$
移到约 1029 $\mathrm{cm^{-1}}$；置换后表面 \ce{DCOO-} 消失而以甲酸蒸气
（\ce{DCOOH}）回到气相 \\
四、毒化／组合 & 用只减少候选物种的预处理，或改变反应物组合，再看活性与
物种量 & 活性随该物种量同比例下降；只有真实反应物组合才生成该物种 &
盐酸或乙酸预处理使 \ce{Al2O3} 表面甲酸离子浓度减到约十分之一而甲酸分解
速率不变（否掉甲酸离子）；甲醇合成中只有 \ce{CO2}$+$\ce{H2} 共吸附才形成
铜甲酸盐 \\
\bottomrule
\end{tabularx}
\end{table}

\begin{figure}[htbp]
\centering
\begin{tikzpicture}
\tikzset{declare function={%
  % 参数名只用单字母，且避开 LaTeX 重音命令 \x、\u（用 \p）。
  % 函数体里的乘法一律加括号，避免 pgfmath 解析歧义。
  ads(\p)=1/(1+((\p-400)/260)*((\p-400)/260));%
  mid(\p)=1/(1+((\p-490)/45)*((\p-490)/45));%
  prd(\p)=1.2/(1+((\p-640)/120)*((\p-640)/120));%
  dcoo(\p)=1/(1+(\p/2.5));%
  hcoo(\p)=1-1/(1+(\p/2.5));%
  gdcoo(\p)=1*(\p/(1+(\p/2)*(\p/2)));}}
\begin{groupplot}[
  group style={group size=2 by 1, horizontal sep=8mm,
               ylabels at=edge left, xlabels at=edge bottom},
  % 尺寸说明：pgfplots 的 width/height 默认**含轴标签与刻度文字**（不是坐标轴盒），
  % 故 54 mm×5.2 cm 实测只得到 37.4 mm×35.4 mm 的轴盒、整图约 96 mm。
  % 这里按 74 mm×7.0 cm 反算，使两列总宽 ≈2×61+8 = 130 mm，贴近单栏图规范
  % 的 \onecolwidth（≈132 mm）而远小于版心 165 mm。
  % 两幅都把 ymax 提到 2.0，是为了给顶左/顶右的图例留出净空——图例按
  % (0.02,0.98)/(0.97,0.98) 锚定在轴顶，其像素高度固定（约 11.5 mm），
  % 提高 ymax 只把曲线整体压低，不移动图例。实测（对产出 PDF 的矢量几何核算）：
  % ymax=1.6 时「候选中间体」峰值与白心圆点标记会擦到图例第三行下沿
  % （重叠 0.2—0.6 mm），ymax=1.4 时面板 (b) 的 HCOO 曲线穿过图例第三行
  % （重叠约 1.3 mm）；改为 2.0 后两处净空均 ≥3 mm。放大尺寸后净空进一步增大。
  width=74mm, height=7.0cm, samples=141,
  axis lines=left,
  x tick label style={/pgf/number format/1000 sep={},
                      font={\fontsize{7}{8.5}\selectfont}},
  y tick label style={font={\fontsize{7}{8.5}\selectfont}},
  xlabel style={font={\fontsize{8}{9.5}\selectfont}},
  ylabel style={font={\fontsize{8}{9.5}\selectfont}},
  ]
  % 公共 domain 会被子面板继承（会报 Dimension too large），故逐面板显式给
  % domain 与 xmin/xmax：面板 (a) 是 400—600 K，面板 (b) 是 0—30 min。
  % 说明：两幅都把 ymax 提到 2.0，是为了给顶左/顶右的图例留出净空——
  % 图例按 (0.02,0.98)/(0.97,0.98) 锚定在轴顶，其像素高度固定（约 11.5 mm），
  % 提高 ymax 只把曲线整体压低，不移动图例。实测（对产出 PDF 的矢量几何核算）：
  % ymax=1.6 时「候选中间体」峰值与白心圆点标记会擦到图例第三行下沿
  % （重叠 0.2—0.6 mm），ymax=1.4 时面板 (b) 的 HCOO 曲线穿过图例第三行
  % （重叠约 1.3 mm）；改为 2.0 后两处净空均 ≥3 mm。
  \nextgroupplot[domain=400:600, xmin=400, xmax=600, ymin=0, ymax=2.0,
                 xtick={400,450,500,550,600}, ytick={0,0.5,1.0,1.5},
                 xlabel={温度 $T$ / K},
                 ylabel={相对强度（模型）},
                 title={(a) 升温序列},
                 title style={font={\fontsize{8}{9.5}\selectfont}},
                 legend style={font={\fontsize{7}{8.5}\selectfont}, draw=none,
                               fill=none, cells={anchor=west},
                               at={(0.02,0.98)}, anchor=north west}]
    % 注意：declare function 的形参只用安全单字母（\p），但 \addplot 的表达式里
    % 必须把绘图自变量 x 作为实参传进去（写成 ads(\p) 会报 Undefined control sequence）。
    \addplot[mainline, black] {ads(x)};
    \addlegendentry{甲酸盐（库）}
    \addplot[thinline, black, dashed, mark=*, mark size=1.1pt,
             mark options={fill=white}] {mid(x)};
    \addlegendentry{候选中间体}
    \addplot[thinline, black, dotted] {prd(x)};
    \addlegendentry{产物速率}
    \draw[gray, dashed] (axis cs:490,0) -- (axis cs:490,1.05);
    \node[anchor=west, font={\fontsize{7}{8.5}\selectfont}]
      at (axis cs:494,0.12) {极大};
  \nextgroupplot[domain=0:30, xmin=0, xmax=30, ymin=0, ymax=2.0,
                 xtick={0,10,20,30}, ytick={0,0.5,1.0,1.5},
                 xlabel={时间 $t$ / min},
                 title={(b) 同位素置换序列},
                 title style={font={\fontsize{8}{9.5}\selectfont}},
                 legend style={font={\fontsize{7}{8.5}\selectfont}, draw=none,
                               fill=none, cells={anchor=west},
                               at={(0.97,0.98)}, anchor=north east}]
    \addplot[mainline, black] {dcoo(x)};
    \addlegendentry{表面 \ce{DCOO-}}
    \addplot[thinline, black, dashed] {gdcoo(x)};
    \addlegendentry{气相 \ce{DCOOH}}
    \addplot[thinline, black, dotted] {hcoo(x)};
    \addlegendentry{表面 \ce{HCOO-}}
\end{groupplot}
\end{tikzpicture}
\caption{把两条判据画成序列：吸附态的「库」与真正中间体在升温过程中的消长
（a），以及同位素置换后表面物种与气相组成的走向（b）}
\label{fig:ch07c-transient}
\modelmark\ 两幅均为模型计算，不是实验谱；曲线用有界洛伦兹型与
$1/(1+t)$ 型函数构造，只用于说明判据的形状特征，不代表任何体系的实际数值。
读者应看的是两处\textbf{形状差异}：(a) 中候选中间体（虚线加圆点）在
$490$\,K 处出现极大，其后被消耗，而产物速率（点线）在更高温度才接近饱和
——「先有极大、后有产物」正是判据一要求的消长关系；储量大的甲酸盐（实线）
则单调缓降，与产物速率没有这种对应。(b) 中置换开始后表面 \ce{DCOO-}（实线）
迅速衰减，同时\emph{气相} \ce{DCOOH}（虚线）出现一个峰、再随吹扫下降，
而表面 \ce{HCOO-}（点线）上升——D 是以反应物分子的形式离开表面回到气相，
而不是以含 D 的产物出现的，这正是判据三、四要区分的情形。局限：(a) 的横轴
是模型温度标尺，(b) 的时间轴按「置换后连续采谱」设定，两者的绝对刻度都不能
用于定量；真实体系中还须叠加气相带扣除与覆盖度校正。
\figsrc{判据（分解速率与总包速率的两个数量级之差、\ce{ZnO} 上两者相等、
DCOOD 的同位素位移与置换后 \ce{DCOO-} 回到气相）据原书 §7.6.1，
书内 p.499—501 / PDF p.518—520；曲线形式为本节模型设定。}
\end{figure}

\begin{measure}
\textbf{反应条件下的原位测量。}以本书给出的原位吸附／反应流程为例
（原书图 7-93）：\ce{H2}、\ce{He}、合成气与乙醇分别经质量流量计计量后进入
原位高温高压漫反射红外原位池，池后接背压调节器再排气，从而在加压条件下
保持反应气氛与压力可控\src{S1 §7.6.3，书内 p.503—504 / PDF p.522—523}。
这套流程的三个要点与判据直接相关：①进料可切换（才能做置换实验）；
②压力与流量可独立设定（才能在同一覆盖度下比分解速率）；③池后可在线的
质谱／色谱取样（才能把表面物种与产物速率放在同一时间轴上）。
\begin{itemize}
  \item \textbf{组合实验}：\ce{Cu/ZnO/Al2O3}（Cu $60\%$、ZnO $30\%$、
        \ce{Al2O3} $10\%$，质量分数）上，只有 \ce{CO2} 与 \ce{H2} 共吸附
        才能形成甲酸盐中间物；多晶 Cu 上甲酸盐分解可给出甲醇，而 \ce{ZnO}
        上甲酸盐分解只给出 \ce{CO2}、CO 与 \ce{H2}
        \src{S1 §7.6.2，书内 p.501—502 / PDF p.520—521}。
  \item \textbf{TPSR 与在线分析}：把 DRIFTS 与 TPSR 结合，同时用在线质谱与
        色谱分析，是阐明吸附物种性质与反应机理最有力的组合。源文强调，
        \ce{Cu-ZnO-Al2O3}（ICI）上甲醇的脱出温度比纯 \ce{ZnO} 上低，
        这才是「甲酸铜盐是甲醇合成关键中间物」的支撑证据（而不是某个带位）
        \src{S1 §7.6.2，书内 p.502—503 / PDF p.521—522}。
  \item \textbf{时间分辨与瞬态}：步进扫描时间分辨红外可把时间窗压到微秒；
        快速扫描 FTIR 用于跟踪 $300$\,ms 光脉冲之后的衰减；水相光催化体系
        用金刚石单晶与 \ce{ZnSe} 晶体组成的全反射池
        \src{S1 §7.6.4，书内 p.510—512 / PDF p.529—531}。
  \item \textbf{电化学条件}：薄层电解池 + 电位差谱 + 锁相检测；参比电位必须
        固定并写明（例如以某一电位的谱作为参考谱），否则「电位依赖」的结论
        无法复现\src{S1 §7.6.4 之后，书内 p.515 / PDF p.534}。
\end{itemize}
\end{measure}

\subsubsection{判例一：甲酸在 \ce{Al2O3} 与 \ce{ZnO} 上的分解}

甲酸分解是催化选择脱氢与脱水的典型反应，也是「吸附态不等于中间体」这一
结论最干净的判例\src{S1 §7.6.1，书内 p.498—501 / PDF p.517—520}。

\begin{obs}
（只列谱学与动力学事实）
\begin{itemize}
  \item[(a)] $100$\,\degC 下甲酸吸附在 \ce{Al2O3} 上，在 $1000$—
        $4000\;\mathrm{cm^{-1}}$ 范围内出现 $2915\;\mathrm{cm^{-1}}$
        （C--H 伸缩）与 $1625\;\mathrm{cm^{-1}}$（O--C--O 反对称伸缩）等带，
        另有 $\delta(\mathrm{C{-}H})$ 与 O--C--O 对称伸缩带；表面存在
        \ce{HCOO-} 吸附态。
  \item[(b)] 改用 \ce{DCOOD} 吸附：O--H 带由 $3580\;\mathrm{cm^{-1}}$ 移到
        约 $2650\;\mathrm{cm^{-1}}$，C--H 带移到约 $2200\;\mathrm{cm^{-1}}$，
        $\delta(\mathrm{C{-}H})$ 移到约 $1029\;\mathrm{cm^{-1}}$；
        谱上没有发现 CO、\ce{CO2} 或分子态甲酸的吸附带。
  \item[(c)] 在反应温度、真空下，\ce{Al2O3} 上 \ce{HCOO-} 吸附态的分解速率
        比同温同覆盖度下甲酸蒸气的分解速率\emph{小两个数量级}。
  \item[(d)] 把 \ce{DCOOD} 分解做到定态后，用 \ce{HCOOH} 迅速置换气相并继续
        反应：表面 \ce{DCOO-} 的带消失，但它以甲酸蒸气（\ce{DCOOH}）的形式
        出现在气相；表面 O--D 迅速被 O--H 置换。
  \item[(e)] 质子交换：\ce{Al2O3} 本身的表面质子可与 \ce{H2O} 中的质子交换；
        甲酸本身的质子在 $180$\,\degC 以下不能交换；而甲酸在 \ce{Al2O3} 上
        解离吸附形成的质子，既可与 \ce{H2O} 中的质子交换，也可与甲酸中的
        质子交换。
  \item[(f)] 用盐酸或乙酸预处理 \ce{Al2O3}，表面甲酸离子浓度降到约十分之一，
        甲酸分解速率不变；吸附的乙酸没有明显分解，也没有观察到乙酸离子与
        甲酸离子之间有明显交换。
  \item[(g)] 由密闭系统的 C、H、O 物料平衡与反应气中反应物、产物分压，得到
        定态总包速率方程，形式为对甲酸一级、被水抑制
        （源文形式：$r=k(\ce{HCOOH})/[1+b(\ce{H2O})]$）。
  \item[(h)] \ce{ZnO} 表面也存在 \ce{HCOO-}，但其行为与 \ce{Al2O3} 不同：
        甲酸在 \ce{ZnO} 上分解的控速步骤是 \ce{HCOO-} 分解，且 \ce{HCOO-}
        的分解速率等于甲酸分解的总包速率。
\end{itemize}
\end{obs}

\begin{infer}
\begin{itemize}
  \item 由 (b) 得：三个带的位移比都落在 $1.3$—$1.4$ 之间（约
        $3580/2650\approx1.35$、$2915/2200\approx1.33$），符合
        式~\eqref{eq:ch07c-isotope} 对含 \ce{X-H} 振动的预言，因此
        $3580\;\mathrm{cm^{-1}}$ 归 O--H、$2915\;\mathrm{cm^{-1}}$ 归 C--H
        的归属被\emph{独立}验证；而 (b) 中不出现分子态甲酸与 CO／\ce{CO2}
        的吸附带，说明该温度下表面物种以甲酸盐与羟基为主。
  \item 由 (c) 与判据二得：\ce{Al2O3} 上的 \ce{HCOO-} 是旁观者，不是中间体。
        这条结论之所以有力，是因为它把「量」与「速率」分开了：\ce{HCOO-}
        的确存在、量也不小，但它的分解慢两个数量级，无法承担总包速率。
  \item 由 (d) 得：\ce{DCOO-} 消失的主要通道是\textbf{被气相置换回甲酸分子}，
        而不是分解成产物——即它是一个可逆吸附的「库」。这一条同时给出了
        判据三的强用法：只看表面物种消失会误判，必须同时看气相组成。
  \item 由 (f) 得：把候选物种的量减少一个数量级而活性不变，直接否掉了
        「甲酸离子是中间体」这一长期默认结论。这是判据四最干脆的一次应用。
  \item 由 (e) 与 (f) 合起来给出一条正向结论：真正参与反应的是甲酸解离吸附
        形成的\textbf{质子}，其作用类似于液相催化中 \ce{H2SO4} 提供质子的
        作用；甲酸根不参与。
  \item 由 (g) 得：速率方程的形式与上面的机理自洽——对甲酸一级（质子的
        来源）、被水抑制（水与甲酸竞争表面位或参与质子转移）。
  \item 由 (h) 得：同一物种在 \ce{ZnO} 上\textbf{是}中间体。因此
        「\ce{HCOO-} 是不是中间体」没有统一答案，必须逐体系按判据二判定。
  \item \textbf{限度}：源文明确指出，这并没有排除在其他反应条件下经甲酸离子
        分解的可能性，只是那条路径需要较高的活化能，要在较高温度下才可能
        发生。判据给出的结论都有条件范围。
\end{itemize}
\end{infer}

\subsubsection{判例二：\ce{Cu/ZnO} 上甲醇合成与乙醇助剂}

\begin{dataflow}
这个体系的信息量最大，也最能说明「多手段合成一条证据链」的做法。链条是：
\textbf{DRIFTS 给物种}（有哪些表面物种、含量随条件怎么变）→
\textbf{TPSR + 在线质谱给温度对应}（哪个物种分解时出现哪个产物）→
\textbf{反应物组合实验给来源}（该物种由哪些反应物生成）→
\textbf{活性数据给结论}（改条件后产物生成速率怎么变）。任何单独一环都不足以
判定中间体：源文列出的三个长期未决问题（C 源是 CO 还是 \ce{CO2}、活性中心
在金属 Cu 表面还是 \ce{ZnO} 表面、Cu 与 \ce{ZnO} 之间是否有协同作用），
正是单靠一种手段回答不了的\src{S1 §7.6.2，书内 p.501 / PDF p.520}。
\end{dataflow}

\textbf{物种与来源。}\ce{Cu/ZnO/Al2O3} 催化剂中 Cu 含量高且高度分散，
透射法不可能用于研究，文献中已发表的甲醇合成红外结果都是在 Cu 含量低于
$15\%$ 时取得的\src{S1 §7.6.2，书内 p.501 / PDF p.520}。改用漫反射后，
本书作者报告：$9\%$ \ce{CO2} $+$ $91\%$ \ce{H2} 于 $200$\,\degC 在 ICI
催化剂上共吸附的 DRIFT 谱与甲醇吸附谱很相似；谱上出现吸附在 Cu 上的表面
甲酸盐的 O--C--O 反对称伸缩带与对称伸缩带，其中对称伸缩带在
$1360\;\mathrm{cm^{-1}}$ 附近；\ce{ZnAl2O4} 上甲酸盐的吸收带、以及单碳酸盐
的吸收带（通常在 \ce{Al2O3} 表面上）位置各不相同\src{S1 §7.6.2，书内
p.502 / PDF p.521}。TPSR 结果给的是判据二与判据四的信息：表面甲酸盐释出时，
气相中出现甲醇以及 \ce{CO2}、\ce{H2}；且 \ce{Cu-ZnO-Al2O3}（ICI）上甲醇的
脱出温度比纯 \ce{ZnO} 上低。由于铜甲酸盐只由 \ce{CO2} 与 \ce{H2} 共吸附
形成，而多晶 Cu 上甲酸盐分解可以给出甲醇、\ce{ZnO} 上甲酸盐分解只给出
\ce{CO2}、CO 与 \ce{H2}，结论是：\textbf{甲酸铜盐是甲醇合成的关键中间物，
而 \ce{ZnO} 上的甲酸盐只是水煤气变换反应的中间物}；\ce{ZnO} 的贡献主要
是通过相互作用调变 Cu 的电子性质（氧阴离子从 \ce{ZnO} 溢流到金属铜上，
增加金属功函数），从而使甲酸盐加氢到甲氧基、再到甲醇更容易进行。

\textbf{乙醇助剂：把中间物换成酯。}工业合成甲醇需高温高压，单程 CO 转化率
只有 $20\%$ 左右；当反应系统中加入乙醇时，反应可在低温低压进行，单程 CO
转化率可达 $50\%$—$80\%$，说明反应路径本身被改变了\src{S1 §7.6.3，
书内 p.503 / PDF p.522}。原位 DRIFTS 给出的表面物种序列是：

\begin{itemize}
  \item 合成气在 \ce{Cu/ZnO} 上先形成甲酸盐物种（\ce{H-COO-Zn} 与
        \ce{H-COO-Cu} 两类）；甲酸盐在 $523$\,K 以下没有被明显加氢，而在
        $523$—$573$\,K 明显分解与加氢，气相中出现甲醇
        \src{S1 §7.6.3，书内 p.505 / PDF p.524}。
  \item 引入乙醇后，谱上出现\textbf{乙氧基}、\textbf{乙酸盐}与
        \textbf{碳酸盐}三类新物种的带；其中碳酸盐与气相乙醇反应生成
        \textbf{甲酸乙酯}物种\src{S1 §7.6.3，书内 p.505—506 / PDF
        p.524—525}。
  \item 甲酸乙酯在石英砂上没有明显的加氢，而在 \ce{Cu/ZnO} 催化剂上被加氢
        成甲醇和乙醇；进一步研究表明，完成加氢的是催化剂中 Cu 上解离吸附的
        氢\src{S1 §7.6.3，书内 p.505 / PDF p.524}。
  \item 由此提出的机理是：合成气（CO、\ce{CO2}、\ce{H2}）在 \ce{Cu/ZnO} 上
        形成甲酸盐物种，与气相或物理吸附的乙醇反应生成甲酸乙酯中间物种，
        最后被 Cu 上的氢还原成甲醇和乙醇（原书图 7-101）。甲酸乙酯是低温
        合成反应的关键中间物，它\emph{改变了经甲氧基的反应机理}
        \src{S1 §7.6.3，书内 p.507—508 / PDF p.526—527}。
\end{itemize}

上述各物种的谱带归属汇总于 \tabref{tab:ch07c-formate}。

\begin{table}[htbp]
\centering
\caption{\ce{Cu/ZnO} 表面相关物种的红外谱带归属\src{据 S1 表 7-19、表 7-20
与图 7-95～图 7-100，书内 p.505—511 / PDF p.524—530 整理。}}
\label{tab:ch07c-formate}
\begin{tabularx}{\onecolwidth}{@{}>{\raggedright\arraybackslash}p{27mm}
  >{\raggedright\arraybackslash}p{33mm} X@{}}
\toprule
表面物种 & 谱带 $/\mathrm{cm^{-1}}$ & 归属与判读要点 \\
\midrule
甲酸盐（Zn 上，\ce{b-HCOO-Zn}） & 2966（sh）、2873；1590（sh）、1572；1370、1362 &
$\nu(\mathrm{C{-}H})$、$\nu_{\mathrm{as}}(\mathrm{OCO})$、$\nu_{\mathrm{s}}(\mathrm{OCO})$
分别对应三组带；三组在表 7-19 中同归 \ce{b-HCOO-Zn} \\
甲酸盐（Cu 上，\ce{b-HCOO-Cu}） & 2925、2849；1542（sh）；1351 &
$\nu(\mathrm{C{-}H})$、$\nu_{\mathrm{as}}(\mathrm{OCO})$、$\nu_{\mathrm{s}}(\mathrm{OCO})$；
铜甲酸盐是甲醇合成的关键中间物（由 \ce{CO2}$+$\ce{H2} 共吸附生成，
TPSR 上甲醇脱出温度与它对应） \\
甲氧基（\ce{CH3O-Zn}） & 2935、2818；1049 &
$\nu(\mathrm{C{-}H})$ 与 $\nu(\mathrm{C{-}O})$；合成气吸附后 ZnO 上即出现
甲氧基，是甲醇合成的另一条路径 \\
表面水（\ce{H2O-Cu}） & 1601 &
$\delta(\mathrm{OH})$；残余表面水的弯曲振动，判读时须与甲酸盐的
$\nu_{\mathrm{as}}(\mathrm{OCO})$ 带区分 \\
乙氧基（乙醇引入后） & 2961、2925、2889、2861；1440、1380；1101、1052 &
$\nu(\mathrm{C{-}H})$、$\delta(\mathrm{CH})$、$\nu(\mathrm{C{-}O})$；
出现乙氧基说明乙醇已解离吸附在表面 \\
羧酸基（乙酸盐，乙醇引入后） & 2933；1562、1451；1049、1023 &
$\nu(\mathrm{C{-}H})$、$\nu(\mathrm{OCO})$ 与 $\rho(\ce{CH3})$；
乙酸物种的出现是乙醇被氧化的旁路证据 \\
碳酸盐（乙醇引入后） & 1551、1534、1428、1345 &
$\nu(\mathrm{COOO})$ 伸缩分裂；碳酸盐可与气相乙醇反应生成甲酸乙酯，
是助剂路径上的关键一环 \\
\bottomrule
\end{tabularx}
\end{table}

\tocheck{未对照原页面图像：原书表 7-19、表 7-20 的列结构在 OCR 中被打乱
（乙氧基／羧酸基／碳酸盐三列的带位互相穿插），上表的带位分组是按源文本可辨
数值与「同一物种的带应成组出现」这一约束重建的；$1360$（甲酸盐对称伸缩）
与源表 $1370/1362$ 几个相近数值的归属、以及 $1551/1534$、$1428/1345$
两组碳酸盐带的配对方式，均未与原页面图像核对。}

\begin{pitfall}
\textbf{把吸附态当中间体。}这是本领域最经典的误判：\ce{Al2O3} 上
\ce{HCOO-} 长期被当成甲酸分解的中间物，直到用判据二测出它的分解速率比总包
速率小两个数量级才被否掉。判据一（序列上的消长）单独用也不够：储量大的物种
在升温过程中常常也「先升后降」（因为总覆盖度下降），必须与产物速率的
\emph{量级}比较。

\textbf{只凭位置归属，不做对照实验。}$1360\;\mathrm{cm^{-1}}$ 这样的带位
可以同时属于甲酸盐、乙酸盐甚至碳酸盐的 O--C--O 伸缩；只有把「共吸附组合
实验」（只有 \ce{CO2}$+$\ce{H2} 才生成）、「温度对应」（TPSR 上产物脱出
温度）与「对照样品」（石英砂上甲酸乙酯不加氢，\ce{Cu/ZnO} 上加氢）放在
一起，归属才成立。

\textbf{忽略气相物种的贡献。}反应条件下的谱图里，气相反应物与产物都有红外
吸收：气相 CO（$2143\;\mathrm{cm^{-1}}$）、气相 \ce{CO2}（$2349\;
\mathrm{cm^{-1}}$）、气相甲酸与乙醇在 $1000$—$3000\;\mathrm{cm^{-1}}$ 有
多条带。若不扣除气相贡献（用背景谱、吹扫对比或差谱），表面物种的带强度与
位置都会被污染；甲酸乙酯这类「气相或物理吸附」的物种更要在图题里写清是
气相带还是表面带。

\textbf{忽略原位池窗口与样品自吸收造成的强度失真。}池窗材料的截止波数、
样品稀释度、装样厚度以及（漫反射下）散射系数的变化，都会改变基线。跨样品
比强度是这一体系上最容易出错的一步；同一体系内部不同温度的强度比较，也要
扣除总覆盖度的变化。

\textbf{把「光照依赖」直接当成光催化中间体的证据。}光照会同时改变温度、
吸附平衡与表面电荷，故光催化体系里必须给出电位或时间依赖，并配合同位素
实验（如 \ce{H2^{18}O} 交换）才能把中间物种与光生载流子过程连起来。
\end{pitfall}

\subsubsection{光催化水氧化：把判据用到光生载流子上}

半导体光催化中，光催化剂先吸收入射光子产生光生电子与空穴，随后经历扩散、
捕获、复合与表面氧化还原反应；水氧化产氧是四空穴转移过程，往往涉及多种
中间物种，特别是不同状态的氧物种，因此原位红外（以及时间分辨红外）在
认识光催化机理中起的作用很重\src{S1 §7.6.4，书内 p.508—510 / PDF
p.527—529}。三条有代表性的证据链：

\begin{itemize}
  \item \textbf{金红石 \ce{TiO2} 上的两个产氧中间物。}用金刚石单晶与
        \ce{ZnSe} 晶体组成的全反射红外池原位研究光激发下金红石 \ce{TiO2}
        纳米颗粒；在 \ce{Fe^{3+}} 作电子牺牲试剂的条件下，光激发时表面出现
        位于 $838$ 与 $812\;\mathrm{cm^{-1}}$ 的两个中间物带，同位素实验将
        它们分别归属为 \ce{TiOOH} 与 \ce{TiOO-Ti}；由此提出 \ce{H2O} 分子
        亲核进攻表面晶格氧空穴生成上述中间物的反应路径
        \src{S1 §7.6.4，书内 p.510 / PDF p.529}。
  \item \textbf{\ce{Co3O4} 上的过氧双核物种。}用快速扫描 FTIR 研究
        \ce{Co3O4} 纳米粒子在 $\mathrm{Ru(bpy)_3^{2+}}$ 光敏剂敏化条件下的
        放氧动力学：$300$\,ms 脉冲光激发后，存在大量未放出的
        $1013\;\mathrm{cm^{-1}}$ 氧中间物种；结合同位素实验与动力学分析，
        归属为三空穴氧化生成的过氧双核物种，并认为是放氧的关键中间物种
        \src{S1 §7.6.4，书内 p.511 / PDF p.530}。
  \item \textbf{直接跟踪光生电子。}光生电子吸收近红外和中红外光，而光生空穴
        在该波段基本没有吸收，不会干扰电子信号，因此时间分辨红外成为研究
        光生电子动力学最有效的技术之一。用步进扫描时间分辨红外研究
        \ce{TiO2}，可观察到中红外范围内的特征宽吸收，其吸光度随波数减小而
        逐渐增加，被归属为浅束缚能级与导带电子；这些电子的衰减在微秒到毫秒
        量级。甲醇存在时，Pt 负载 \ce{TiO2} 上光生电子在微秒量级基本不衰减，
        直到毫秒到秒才明显衰减，增加气相甲醇或水汽可明显促进这一衰减；结合
        光催化重整甲醇制氢的活性结果，提出分子态吸附的水与甲醇促进了质子
        转移，从而促进 Pt 上的产氢反应
        \src{S1 §7.6.4，书内 p.512—513 / PDF p.531—532}。
\end{itemize}

\tocheck{未对照原页面图像：金红石 \ce{TiO2} 上两个中间物带位的文字在 OCR 中
可辨为 838 与 $812\;\mathrm{cm^{-1}}$，但配图图注（图 7-102）中的第二处带位
与归属箭头无法辨读；时间分辨一段中样品的标号（源文作「P?25TiO2」）不能确定
是 Pt/TiO2 还是另一种 Pt 负载量标号，本分片按「Pt 负载 TiO2」表述。}

\subsubsection{电化学与光电催化原位红外：薄层、差谱与毒性中间体}

\emph{（本小节标题及其前一段文字所在的页面在源文本中缺损，标题为按内容拟定。）}
\tocheck{未对照原页面图像：书内 p.514（PDF p.533）的 OCR 层为空，该页上的
小节标题与部分正文缺失；本小节标题「电化学与光电催化原位红外」是按源文
内容（固液界面红外、薄层电解池、电催化机理）拟定的，未与原页面核对。}

固液界面的红外研究要克服三个障碍：①溶剂分子（通常为水）对红外能量的大量
吸收；②固体电催化剂表面反射红外光导致部分能量损失；③表面吸附分子量少，
满单层情况下吸附量仅在 $10^{-9}$\,mol\,cm$^{-2}$ 量级。这些障碍使信号非常
微弱，需用薄层电解池、电位差谱与各种微弱信号检测技术提高信噪比
\src{S1 §7.6.4 之后，书内 p.515 / PDF p.534}。常用的做法是让工作电极与
红外窗片紧密接触，形成很薄的液体薄层电解池，以尽量减少红外光被溶剂吸收的
量；液层厚度是这套方法的关键参数（源文引用的研究表明，只有当液层厚度小到
一定程度时才能得到较高信噪比）。电位差谱则是在保持其他实验条件不变的情况下
仅改变电极电位，再做差减归一化；锁相检测、偏振调整等技术也常用于微弱信号
检测。\tocheck{未对照原页面图像：液层厚度上限的数值在 OCR 中作「130」后接
一个损坏字符（可能为 µm），本分片未写入具体上限，只保留「需小到足够薄」
这一条件；满单层吸附量的数量级（$10^{-9}$\,mol\,cm$^{-2}$）在 OCR 中同样
残缺，按学科常见量级补写，未与原文核对。}

判据在这里的用法与前面一致，只是「序列」的变量换成了\textbf{电位}：

\begin{itemize}
  \item \textbf{CO 作探针。}在经过电化学处理的多晶 Rh 电极上可检测到两种
        CO 孪生吸附态：用较快的电位扫描速度（$1.5$\,V/s）在 $2.4$\,V 电位
        区间处理电极，可在表面形成一层 Rh 原子簇氧化物膜，使吸附 CO 给出一对
        宽的 IR 吸收带；当表面金属氧化物被还原后，CO 仍以孪生态吸附在电极
        表面，但吸收带红移到 $2012$ 与 $2032\;\mathrm{cm^{-1}}$，半峰宽减小
        \src{S1 §7.6.4 之后，书内 p.516 / PDF p.535}。
  \item \textbf{毒性中间体与结构敏感性。}甲醇的电催化氧化遵循双途径机理
        （活性中间体途径与经毒性中间体的途径）。电化学原位红外在 Pt 单晶
        电极上区分出两类毒性中间体：位于 $2050\;\mathrm{cm^{-1}}$ 左右的
        线式吸附 CO，以及多重态吸附物种；毒性中间体的起始氧化电位顺序为
        Pt(111) $<$ Pt(110) $<$ Pt(100)，说明甲醇电催化氧化是表面结构敏感
        反应\src{S1 §7.6.4 之后，书内 p.517 / PDF p.536}。
  \item \textbf{助剂如何减少毒性中间体。}在 Pd/C 电极上可以观测到 CO 带，
        而硼的引入使 Pd 表面积累的 CO 发生实质性减少；这些吸附 CO 应是主要
        的毒性中间体，其减少显著改善了甲酸燃料电池的性能（原书图 7-108 以
        $1.25$\,V 的谱作参考谱）\src{S1 §7.6.4 之后，书内 p.517 / PDF
        p.536}。
  \item \textbf{光电催化水氧化。}在 \ce{Fe2O3} 光阳极上，于光电催化水氧化
        条件下观测到 $898\;\mathrm{cm^{-1}}$ 处的新红外吸收中间物种；该吸收
        对光照存在明显依赖，仅在光电催化水氧化发生时存在，被归属为高价铁氧
        中间物种 $\mathrm{Fe^{IV}{=}O}$，即水氧化第一步空穴氧化过程的中间物
        \src{S1 §7.6.4 之后，书内 p.517—518 / PDF p.536—537}。
\end{itemize}

\begin{crosscheck}
\textbf{与电化学曲线互查。}中间物种的电位依赖必须与 $J$--$V$ 曲线、
起始氧化电位、库仑量对应；若某个带只随电位移动而与电流无关，更可能是双电层
或溶剂重组效应。

\textbf{与同位素、质谱互查。}光催化与电催化体系里，\ce{H2^{18}O} 交换可以
把「含水物种」与「含晶格氧物种」分开；微分电化学质谱可以给出产物的实时
生成速率，与红外带强度同轴比较，正是判据二在电化学条件下的形式。

\textbf{与理论计算互查。}带位归属（如 \ce{TiOOH} 与 \ce{TiOO-Ti} 的区分、
$\mathrm{Fe^{IV}{=}O}$ 的指认）最终要与频率计算和反应能垒计算一致；计算不能
替代判据，但可以排除几张谱之间无法区分的候选结构。

\textbf{回到方法边界。}红外给出的是有红外活性的物种；对称性高、偶极变化小的
物种（以及被强吸收带掩盖的物种）可能\emph{测不到}。「没看到」永远不能直接
推出「不存在」，这也是本书在光催化与电催化两节反复强调的边界。
\end{crosscheck}

\srcnote{本分片来源：S1（上册）§7.5—§7.6，以及 §7.6 之后的电化学与光电催化
原位红外部分，书内 p.489—518 / PDF p.508—537。其中书内 p.514（PDF p.533）的
OCR 层缺损，该处小节标题与部分正文缺失，已在正文用 \tocheck{} 标注。
原书参考文献与作者简介见书内 p.518—523 / PDF p.537—542，本分片不复述。
本节内所有谱带数值均取自源文本文字层，凡 OCR 残缺或互相矛盾之处均已在正文
就地标注，未对照原页面图像。}
